% Matthieu Greiner, resume (English, US Letter)
% Build: pdflatex cv_en.tex
\documentclass[10pt]{article}
\usepackage[letterpaper,top=0.45in,bottom=0.4in,left=0.6in,right=0.6in]{geometry}
\usepackage[T1]{fontenc}
\usepackage[utf8]{inputenc}
\usepackage[default]{sourcesanspro}
\usepackage{fontawesome5}
\usepackage{xcolor}
\usepackage{titlesec}
\usepackage{enumitem}
\usepackage{microtype}
\usepackage[hidelinks]{hyperref}

\definecolor{muted}{HTML}{4A4A4A}
\pagestyle{empty}
\setlength{\parindent}{0pt}
\setlength{\parskip}{0pt}
\linespread{1.0}

\titleformat{\section}{\scshape\large}{}{0pt}{}[\vspace{-2pt}\titlerule]
\titlespacing*{\section}{0pt}{10pt}{5pt}

\setlist[itemize]{leftmargin=1.2em,itemsep=0pt,topsep=1pt,parsep=0pt,partopsep=0pt,label={\textendash}}

\newcommand{\role}[4]{\textbf{#1} \textbar{} \textit{#2}\hfill\textit{#3}\par\vspace{0pt}}
\newcommand{\school}[4]{\textbf{#1}\hfill\textit{#2}\par\textit{#3}\hfill\textit{#4}\par}
\newcommand{\gap}{\vspace{4pt}}

\begin{document}
\fontsize{10.4}{12.4}\selectfont

{\LARGE\bfseries Matthieu Greiner}\par\vspace{2pt}
{\small\itshape
\faMapMarker* Paris, France \,|\,
\faEnvelope \ \href{mailto:greinermatthieu635@gmail.com}{greinermatthieu635@gmail.com}\par
\faGlobe \ \href{https://matthieugreiner.com}{matthieugreiner.com} \,|\,
\faLinkedin \ \href{https://www.linkedin.com/in/matthieu-greiner}{linkedin.com/in/matthieu-greiner} \,|\,
\faGithub \ \href{https://github.com/Spaghetto784}{github.com/Spaghetto784}\par}
\vspace{5pt}
{\small Cybersecurity engineering student at EPITA (class of 2028), DevSecOps apprentice and malware researcher.
Seeking a 6-month internship in the United States starting early 2028.\par}

\section{Education}
\school{EPITA, Engineering School of Computer Science}{Paris}{Engineering degree (Master's level) in Computer Science, Cybersecurity major}{2023 -- 2028 (expected)}\gap
\school{Chulalongkorn University}{Bangkok}{Exchange semester, Computer Engineering}{2025}\gap
\school{Universit\'e Paris-Est Cr\'eteil (UPEC)}{Paris}{Bachelor in Computer Science}{2023 -- 2026}

\section{Experience}
\role{DevSecOps Engineer (Apprentice)}{intuitem, Vélizy}{Sep 2026 -- Present}{}
\begin{itemize}
\item Build features for CISO Assistant, an open-source GRC platform (Django, Svelte; 4,000+ GitHub stars, 3M+ downloads); \href{https://github.com/intuitem/ciso-assistant-community/pulls?q=is:pr+is:merged+author:Spaghetto784}{merged pull requests} include data-wizard import/export, custom field types and EBIOS RM fixes
\item Keep every change compliant with the CI/CD pipeline; work with Kubernetes and multi-cloud environments
\item Answer customer requests on the Jira support desk
\end{itemize}\gap
\role{Student Researcher}{EPITA Research Lab (LRE)}{Oct 2026 -- Present}{}
\begin{itemize}
\item Study malware evasion techniques on recent Windows x64 samples with dynamic binary instrumentation (Intel Pin)
\item Track how evasion techniques evolve over time; goal: publish a paper before graduation (2028)
\end{itemize}\gap
\role{Information Systems \& Security Engineer (Apprentice)}{SFR Business, Paris}{2025 -- 2026}{}
\begin{itemize}
\item Supported security governance and compliance (NIS2, DORA, GDPR, ISO 27001); ran risk analyses (EBIOS RM), audits and remediation tracking
\item Built CyberTool, an internal Go utility checking TLS, DNS, HTTP responses and security headers, used by a 12-person team
\item Contributed to an internal AI assistant for security questionnaires, adopted by the team and boosting its KPIs; automated security KPI reporting (Excel/VBA)
\end{itemize}\gap
\role{Software Development Intern}{Entrepose \'Echafaudages}{2025}{}
\begin{itemize}\item Developed a C\# plugin for Autodesk Robot and automated engineering workflows\end{itemize}

\section{Projects}
\textbf{SecureFlow} \textbar{} \textit{DevSecOps platform}\hfill\textit{github.com/Spaghetto784/SecureFlow}\par
React, FastAPI, PostgreSQL, Docker; CI security scans (Semgrep, Gitleaks, Trivy); JWT and Argon2 auth.\gap\par
\textbf{wlkom-rootkit} \textbar{} \textit{EPITA systems project}\hfill\textit{C, Python}\par
Educational Linux kernel module built in an isolated two-VM QEMU lab, with a secure TCP client/server protocol.\gap\par
\textbf{Budgify} \textbar{} \textit{Privacy-first iOS app}\hfill\textit{Swift, SwiftUI}\par
On-device ML categorization, Face ID, AES-GCM encryption and Keychain key management; no external services.

\section{Technical Skills}
\textbf{Languages}: C, Go, Python, Java, C\#, Swift, SQL, JavaScript, TypeScript\par
\textbf{Systems}: Linux, kernel modules, QEMU/KVM, GDB\par
\textbf{Security}: malware analysis, Intel Pin (DBI), vulnerability analysis, TLS/DNS\par
\textbf{DevOps \& Cloud}: Git, Docker, Kubernetes, CI/CD, REST APIs, AWS, Azure\par
\textbf{Frameworks}: Django, Svelte, FastAPI, React\par
\textbf{GRC}: EBIOS RM (certified Risk Manager), ISO 27001, NIS2, DORA, GDPR; CompTIA Security+ (in progress)

\section{Languages \& Interests}
\textbf{Languages}: French (native), Armenian (native), English (Cambridge C1), Spanish (B1)\par
\textbf{Interests}: electric guitar, CTF competitions, reading

\end{document}
